\section{Prime-power sums and zeros}\label{sec:analytic}

\subsection{The Laplace transform of the remainder}

Let \(d\) be an odd prime and \(q\) a positive integer coprime to \(d\).
Write \(\chi_d(n)=(n/d)\) for the Legendre symbol, \(\Lambda\) for the von Mangoldt function,
and \(\varphi\) for Euler's totient function. Put
\[
  \lambda_q(n)=\mathbf 1_{n\equiv1\bmod q}\Lambda(n)(1-\chi_d(n)).
\]
For the moment, \(d\) is fixed and is suppressed from the notation \(\lambda_q,R_q,G_q\).
If \(\psi\) is principal, let \(L^\circ(s,\psi)=(s-1)L(s,\psi)\), with the value at
\(s=1\) given by holomorphic continuation. For nonprincipal \(\psi\), put
\(L^\circ(s,\psi)=L(s,\psi)\). Both definitions give entire functions.
We regard \(\psi\chi_d\) as a character modulo \(qd\). Since \(d\) is an odd prime
and \((d,q)=1\), \(\psi\chi_d\) is nonprincipal. Set
\begin{align}
  G_q(s)
   &=\frac1{\varphi(q)}
     \left[
       \sum_{\psi\bmod q}
       \left(
        \frac{(L^\circ)'(s,\psi\chi_d)}{L^\circ(s,\psi\chi_d)}
        -\frac{(L^\circ)'(s,\psi)}{L^\circ(s,\psi)}
       \right)-1
     \right], \label{eq:log-derivative}\\
  R_q(y)
   &= \sum_{n\le y}\lambda_q(n)\left(\frac1n-\frac1y\right)
      -\frac{\log y}{\varphi(q)}-c_q,\qquad c_q=\RePart G_q(1).
  \label{eq:prime-remainder}
\end{align}
Logarithmic derivatives are taken with respect to \(s\). For a principal character,
\(L^\circ\) is nonzero at \(s=1\) after removal of the pole. For a nonprincipal
character, \(L(1,\psi)\ne0\). Thus \(G_q\) is holomorphic at \(1\).

By character orthogonality, \(\lambda_q(n)\) is the average of
\(\Lambda(n)(\psi(n)-(\psi\chi_d)(n))\) over all characters modulo \(q\).
Logarithmic differentiation of the Euler products for \(\RePart s>1\), with the term
\(1/(s-1)\) from removing the pole of the principal \(L\)-function accounted for, gives
\begin{equation}\label{eq:lambda-dirichlet}
  F_q(s):=\sum_{n\ge1}\frac{\lambda_q(n)}{n^s}
        =\frac{s}{\varphi(q)(s-1)}+G_q(s).
\end{equation}
The right-hand side is real for real \(s>1\), so \(G_q(1)=c_q\).

\begin{lemma}\label{lem:prime-transform}
For \(\RePart s>2\) one has
\begin{equation}\label{eq:prime-transform}
  \mathcal T_q(s):=\int_0^\infty R_q(e^v)e^{(1-s)v}\,dv
     =\frac{G_q(s)/s-G_q(1)}{s-1}.
\end{equation}
The meromorphic continuation given by the right-hand side is holomorphic at \(s=1\).
At a point \(\rho\ne1\) with \(\RePart\rho>0\), it has at most a simple pole, with residue
\[
  \res_{s=\rho}\mathcal T_q(s)=
    \frac{\nu(\rho)-\nu_d(\rho)}{\varphi(q)\rho(1-\rho)}.
\]
Here \(\nu\) and \(\nu_d\) are the sums over all characters modulo \(q\) of the
orders of vanishing at \(\rho\) of \(L^\circ(s,\psi)\) and \(L^\circ(s,\psi\chi_d)\),
respectively.
\end{lemma}

\begin{proof}
Put \(A_r(y)=y^{-r}\sum_{n\le y}\lambda_q(n)n^{r-1}\).
For \(r=0,1\), integrating the \(n\)-th summand over \(v\ge\log n\) gives
\[
 \int_0^\infty A_r(e^v)e^{(1-s)v}\,dv=\frac{F_q(s)}{s+r-1}.
\]
The termwise integration is absolutely convergent for \(\RePart s>2\). Consequently
\[
 \mathcal T_q(s)
 = F_q(s)\left(\frac1{s-1}-\frac1s\right)
    -\frac1{\varphi(q)(s-1)^2}-\frac{G_q(1)}{s-1}.
\]
Substitution of \eqref{eq:lambda-dirichlet} gives \eqref{eq:prime-transform}.
The numerator vanishes at \(s=1\). At \(s=\rho\ne1\), apply to
\eqref{eq:log-derivative} the fact that the residue of the logarithmic derivative
of an entire function equals its order of vanishing.
\leanlinks{\leanref{integral_primeRemainder}}
\end{proof}

\subsection{Selecting the quadratic twist}

\begin{lemma}\label{lem:select-quadratic}
Let \(q\) be a positive multiple of \(8\), with \(d_0\equiv7\pmod8\) and \((d_0,q)=1\).
For every \(\beta>1/2\) and \(H\ge0\) there is a prime \(d\equiv d_0\pmod q\)
for which \(L^\circ(s,\psi\chi_d)\) is nonzero throughout
\(\RePart s\ge\beta,\ |\ImPart s|\le H\), for every character \(\psi\) modulo \(q\).
\end{lemma}

\begin{proof}
For \(\beta\ge1\), the case \(\beta=3/4\) suffices.
Assume \(1/2<\beta<1\) and write \(\beta=1/2+9\eta\), with \(\eta>0\).
The formal proof gives the bound \(O_{q,\eta,H}(X^{1-\eta/2})\) for the number
of primes \(d\le X\) for which at least one twisted function has a zero
in the specified region, by the following argument.

Let \(\mu\) be the Möbius function. For \(j\in\{0,1\}\) put
\(w_j(u)=u^j e^{-\pi u^2}\) and write
\(\Gamma_j(s)=\int_0^\infty w_j(u)u^{s-1}\,du\) for its Mellin transform.
For a sequence of coefficients \(c=(c(n))_{n\ge1}\), set
\[
 \theta_c(u)=u^{1/2}\sum_{n\ge1}c(n)w_j(nu),\qquad
 J(c,T)=\int_{1/T}^\infty
        \bigl(|\theta_c(u)|+|\theta_{\mu c}(u)|\bigr)\,\frac{du}{u}.
\]
Let \(\chi\) be a nonprincipal primitive character of conductor \(Q\le T^2\)
and parity \(j\). The Mellin transform of \(\theta_\chi\) at \(s-1/2\) equals
\(\Gamma_j(s)L(s,\chi)\). For \(\RePart s>1\), the corresponding integral for
\(\theta_{\mu\chi}\) gives \(\Gamma_j(s)/L(s,\chi)\). Let \(M_T(s)\) denote
that integral restricted to \(u>1/T\). On \(\RePart s=1/2\) we have
\(|M_T(s)|\le J(\chi,T)\).
The functional equation for the theta function gives
\(\int_0^\infty |\theta_\chi(u)|\,du/u\) as twice
\(\int_{Q^{-1/2}}^\infty|\theta_\chi(u)|\,du/u\).
Since \(Q^{-1/2}\ge T^{-1}\), it follows that
\(|\Gamma_j(s)L(s,\chi)|\le2J(\chi,T)\) on this line.

Suppose \(J(\chi,T)\le T^{2\eta}\) and put
\(\Psi(s)=\Gamma_j(s)L(s,\chi)M_T(s)-\Gamma_j(s)^2\).
On the left boundary \(\RePart s=1/2\) we have \(|\Psi(s)|\ll T^{4\eta}\).
On the right boundary \(\RePart s=2\), the bound \(3e^{-u}\) for the
Mellin integrand of a coefficient sequence bounded by \(1\) in modulus gives
\[
 |\Gamma_j(s)L(s,\chi)|\le3,\qquad
 \left|\frac{\Gamma_j(s)}{L(s,\chi)}-M_T(s)\right|\le 3/T.
\]
Thus \(|\Psi(s)|\le9/T\) there. The Mellin representation of \(\Gamma_j L\) and
the truncated integral defining \(M_T\) give the holomorphy and boundedness
required by the three-lines theorem on the strip. For
\(\RePart s\in[1/2+9\eta,1]\), the interpolation weight of the right boundary is
\(t=(\RePart s-1/2)/(3/2)\ge6\eta\).
Hence \(|\Psi(s)|\ll T^{4\eta(1-t)-t}\), and the exponent is at most \(-2\eta\).
The Euler product rules out zeros in \(\RePart s>1\).
If \(L(s,\chi)=0\), then \(\Psi(s)=-\Gamma_j(s)^2\).
The modulus of \(\Gamma_j\) has a positive lower bound on
\(1/2\le\RePart s\le1,\ |\ImPart s|\le H\). Therefore a zero in the required
region forces \(J(\chi,T)>T^{2\eta}\) for all sufficiently large \(T\).
\leanlinks{\leanref{eventually_rpow_lt_thetaTailIntegral}}

Next let \(n^\circ\) be the odd part of \(n\), and for \(|\alpha_n|\le1\) put
\(c_v(n)=\alpha_n(-v/n^\circ)\), using the Jacobi symbol.
When \(u\ge T^{-1}\), expansion of \(\sum_{1\le v\le T^2}|\theta_{c_v}(u)|^2\)
gives the character sum
\[
  \sum_{1\le v\le T^2}\left(\frac{-v}{(nn')^\circ}\right)
\]
for the pair \(n,n'\). If the odd part of \(nn'\) is a square, the absolute
value of the sum is at most \(T^2\). With \(A=1/2+\eta\) in
\(|w_j(x)|\ll_A x^{-A}\), the contribution of the pair is bounded by a constant
times \(T^2u^{-2\eta}(nn')^{-1/2-\eta}\).
For such a pair there is a \(k\in\{1,2\}\) for which \(knn'\) is a square, so
we may write \(kn=ab^2,\ n'=ac^2\). The sum \(\sum(nn')^{-1/2-\eta}\) over these
pairs is therefore bounded by a constant times the convergent triple sum
\(\sum_{a,b,c\ge1}(abc)^{-1-2\eta}\).

If the odd part is not a square, the Pólya--Vinogradov inequality for the Jacobi
character gives \(O_\eta((nn')^{1/2+\eta/2})\) for the character sum, after
absorbing the logarithmic factor. Taking \(A=3/2+\eta\) and using \(u^{-2}\le T^2\)
bounds the contribution of a pair by a constant times
\(T^2u^{-2\eta}(nn')^{-1-\eta/2}\). This double sum also converges. Combining the
two bounds gives
\[
 \sum_{v\le T^2}|\theta_{c_v}(u)|^2\ll_\eta T^2u^{-2\eta}.
\]
Cauchy--Schwarz followed by integration against \(du/u\) yields
\(\sum_{v\le T^2}J(c_v,T)\ll_\eta T^{2+\eta}\).
For the term involving \(\theta_{\mu c_v}\), the coefficients are
\(\alpha_n\mu(n)\).
\leanlinks{\leanref{sum_thetaTailIntegral_le}}

For \(d\equiv7\pmod8\), quadratic reciprocity and \((2/d)=1\) give
\((n/d)=(-d/n^\circ)\). For each \(\psi\bmod q\), let \(\psi^\ast\) be the
primitive character inducing \(\psi\), with conductor \(Q_0\mid q\).
Then \(\psi^\ast\chi_d\) is primitive and nonprincipal, with conductor \(Q_0d\).
The finitely many Euler factors introduced by induction are nonzero in
\(\RePart s>0\). Set \(T^2=qX\) and \(\alpha_n=\psi^\ast(n)\).
Counting each parity separately, the number of primes producing a zero
in the region is \(O(T^{2+\eta}/T^{2\eta})=O(X^{1-\eta/2})\).
There are finitely many characters modulo \(q\), so the bound has the same
exponent after summing over them.

Were every prime in the residue class to produce such a zero, this counting
estimate would imply convergence of \(\sum_{p\equiv d_0\bmod q}(\log p)/p\).
That series diverges in a reduced residue class.
\leanlinks{\leanref{zero_density_quadTwist}}
\end{proof}

\subsection{Zero exclusion from a one-sided bound}

\begin{lemma}[Landau's lemma in the required form]\label{lem:landau}
Let \(U(v)\) be a measurable real function satisfying \(|U(v)|\le C e^v\)
for \(v\ge0\). Suppose its transform \(T(s)=\int_0^\infty U(v)e^{(1-s)v}\,dv\)
has a meromorphic continuation to \(\RePart s>0\) with at most simple poles.
Suppose also that, for some \(\beta>0\), the inequality \(U(v)\le e^{(\beta-1)v}\)
holds for all sufficiently large \(v\), and the real part of the residue of
\(T\) at each real \(\sigma>\beta\) is nonnegative. Then \(T\) is holomorphic
on \(\RePart s>\beta\).
\end{lemma}

\begin{proof}
Choose \(K\) so that \(K e^{\beta v}-U(v)e^v\ge0\) for all \(v\ge0\).
The Laplace transform \(A(s)\) of this nonnegative density satisfies
\(A(s)=K/(s-\beta)-T(s)\) sufficiently far to the right.
Let \(\sigma_c\) be its abscissa of convergence and suppose \(\sigma_c>\beta\).
If \(A\) were holomorphic at \(\sigma_c\), then for \(\sigma_1>\sigma_c\)
sufficiently close to \(\sigma_c\) the Taylor series about \(\sigma_1\) would
converge at some real \(\sigma_2<\sigma_c\). Differentiation under the integral
in the half-plane of convergence identifies its \(j\)-th term with the integral
of the density against \(e^{-\sigma_1 v}((\sigma_1-\sigma_2)v)^j/j!\).
All terms are nonnegative, so monotone convergence makes the sum equal to the
Laplace integral at \(\sigma_2\), and that integral is finite. This contradicts
the definition of the abscissa of convergence.

Thus \(A\) has a pole at \(\sigma_c\). Its residue is the negative of the residue
of \(T\) there, so its real part is nonpositive. At a simple pole, however, the
residue is the limit
\(\lim_{\sigma\downarrow\sigma_c}(\sigma-\sigma_c)A(\sigma)\).
This is a nonnegative real number and is nonzero at a simple pole, hence is
strictly positive. The contradiction gives \(\sigma_c\le\beta\).
The integral, and therefore \(K/(s-\beta)-A(s)\), is holomorphic on
\(\RePart s>\beta\).
\leanlinks{\leanref{IsPolar.landau}}
\end{proof}

\begin{proposition}\label{prop:one-sided-criterion}
Let \(q,d_0\) be as in Lemma \ref{lem:select-quadratic} and let \(\beta>1/2\).
Suppose that for every prime \(d\equiv d_0\pmod q\) one has
\(R_q(y)\le y^{\beta-1}\) for all sufficiently large real \(y\).
Then \(L(\rho,\psi)\ne0\) for every \(\psi\bmod q\) and every \(\rho\ne1\)
with \(\RePart\rho>\beta\).
\end{proposition}

\begin{proof}
Fix \(\rho\) and choose \(d\) by Lemma \ref{lem:select-quadratic}, with
\(H=|\ImPart\rho|\). The residue in Lemma \ref{lem:prime-transform} is then
\(\nu(s)/\bigl(\varphi(q)s(1-s)\bigr)\) for
\(\RePart s>\beta,\ |\ImPart s|\le H,\ s\ne1\).
In particular, it is nonnegative on the real interval \(\beta<s<1\).
For real \(s\ge1\), the entire functions \(L^\circ(s,\psi)\) are nonzero and the
residue vanishes. The function \(U(v)=R_q(e^v)\) is \(O(e^v)\), and the
hypothesis gives \(U(v)\le e^{(\beta-1)v}\) for all sufficiently large \(v\).
By Lemma \ref{lem:landau} the residue at \(\rho\) is zero.
The denominator is nonzero, and \(\nu(\rho)\) is a sum of nonnegative orders
of vanishing, so every summand is zero. Since \(\rho\ne1\), nonvanishing of
\(L^\circ\) implies nonvanishing of \(L\).
\leanlinks{\leanref{LFunction_ne_zero_of_primeRemainder_le}}
\end{proof}

\subsection{The prime number theorem for a fixed modulus}

\begin{lemma}[Uniform logarithmic derivative bound]\label{lem:uniform-log-derivative}
Let \(\chi\) be a Dirichlet character modulo \(Q\), with \(\mathcal L\ge10^{10}\)
and \(Q\le e^{\mathcal L}\). Write \(\delta_\chi=1\) for the principal character and
\(\delta_\chi=0\) otherwise, and put \(f_\chi(s)=L^\circ(s,\chi)/s^{\delta_\chi}\).
If \(\chi^2\) is principal, assume that \(|f_\chi(z)|\ge1/\mathcal L\) for
\(|z-1|\le1/\mathcal L\). Then for \(\RePart s\ge1-\mathcal L^{-25}\) and
\(|s|\le e^{\mathcal L}\) one has
\begin{equation}\label{eq:uniform-log-derivative}
  L^\circ(s,\chi)\ne0,\qquad
  \left|\frac{(L^\circ)'(s,\chi)}{L^\circ(s,\chi)}\right|\le2\mathcal L^{25}.
\end{equation}
\end{lemma}

\begin{proof}
Put \(\kappa_\chi=Q^{-1}\sum_{v\bmod Q}\chi(v)\) and
\(r_\chi(t)=\sum_{1\le n\le t}\chi(n)-\kappa_\chi t\).
Subtracting the contribution of each complete period leaves fewer than \(Q\)
terms, each of absolute value at most \(2\), and a contribution of at most \(1\)
from \(t-\lfloor t\rfloor\). Thus \(|r_\chi(t)|\le2Q\).
Partial summation and analytic continuation give
\[
 f_\chi(s)=\kappa_\chi+(s-\delta_\chi)
            \int_1^\infty r_\chi(t)t^{-s-1}\,dt
 \qquad(\RePart s>0).
\]
Indeed, in \(\RePart s>1\) the \(L\)-function is \(s\) times the integral
plus \(s\kappa_\chi/(s-1)\). For a nonprincipal character \(\kappa_\chi=0\).
For the principal character the factor \((s-1)/s\) removes the pole.
Since \(|s-\delta_\chi|\le|s|\) in \(\RePart s\ge1/2\), the integral gives
\(|f_\chi(s)|\le1+4Q|s|\) there.

Fix \(s\) with \(\RePart s\ge1-1/\mathcal L\) and \(|s|\le3e^{\mathcal L}\).
On the vertical strip of width \(1/3\) with right edge
\(\RePart s+2/\mathcal L\), consider \(g(z)=f_\chi(z)/(z-s+2)\).
Throughout the strip \(|z-s+2|\ge1\) and \(|z|\le4e^{\mathcal L}|z-s+2|\),
so the partial summation bound gives \(|g(z)|\le e^{3\mathcal L}\).
The right edge is in \(\RePart z\ge1+1/\mathcal L\), where bounding the
Dirichlet series by \(\zeta(\RePart z)\) gives \(|g(z)|\le1+\mathcal L\).
The weight of the left edge in the three-lines theorem is \(6/\mathcal L\).
Consequently
\[
 |f_\chi(s)|=2|g(s)|\le2e^{18}(1+\mathcal L)\le\mathcal L^2.
\]
For \(\RePart z\ge 1-1/(2\mathcal L)\) and \(|z|\le2e^\mathcal L\), applying
this bound on the circle of radius \(1/(2\mathcal L)\) gives
\(|f_\chi'(z)|\le2\mathcal L^3\) by Cauchy's estimate.
\leanlinks{\leansrc{Lib/Analytic/PrimeNumberTheorem/LFunctionGrowth}{Partial summation, three lines, and the derivative bound}}

Now set \(\eta=\mathcal L^{-25}\) and take \(v=\sigma+it\) with
\(\sigma\ge1+\eta\) and \(|v|\le3e^\mathcal L/2\).
In the logarithms of the Euler products, the coefficient at each prime
power with \(p\nmid Q\) is multiplied by
\(3+4\cos\theta+\cos(2\theta)=2(1+\cos\theta)^2\ge0\), where \(\theta\) is
the argument of \(\chi(p)^j p^{-ijt}\). Hence
\[
 1\le |L(\sigma,1_Q)|^3 |L(v,\chi)|^4 |L(\sigma+2it,\chi^2)| .
\]
The first \(L\)-value is bounded by \(2/\eta\) using the zeta function in
the region of absolute convergence. If \(\chi^2\ne1_Q\) or
\(|v-1|\ge1/\mathcal L\), the other two \(L\)-values are bounded by
\(2\mathcal L|f_\chi(v)|\) and \(2\mathcal L^3\), respectively.
For a principal character the factor that is restored has modulus
\(|v|/|v-1|\le1+\mathcal L\).
The same estimate applies at the second point since
\(|\sigma+2it-1|\ge|v-1|\) and \(|\sigma+2it|\le3e^\mathcal L\).
Thus \(\eta^3\le256\mathcal L^7|f_\chi(v)|^4\). From
\(256\cdot12^4\le\mathcal L^6\) we obtain \(|f_\chi(v)|\ge12/\mathcal L^{22}\).
When \(\chi^2=1_Q\) and \(|v-1|<1/\mathcal L\), the hypothesis yields
the same bound since \(1/\mathcal L\ge12/\mathcal L^{22}\).

For \(s\) in the region of the lemma, the lower bound applies at \(v=s+2\eta\).
The derivative is bounded by \(2\mathcal L^3\) along the segment \([s,v]\).
It follows that
\[
 |f_\chi(v)-f_\chi(s)|\le4/\mathcal L^{22},
 \qquad |f_\chi(s)|\ge8/\mathcal L^{22}.
\]
Therefore \(|f_\chi'(s)/f_\chi(s)|\le\mathcal L^{25}/4\).
On passing back to \(L^\circ(s,\chi)=s^{\delta_\chi}f_\chi(s)\), the added
term in the logarithmic derivative is \(\delta_\chi/s\), of modulus at
most \(2\). This proves \eqref{eq:uniform-log-derivative}.
\leanlinks{\leansrc{Lib/Analytic/PrimeNumberTheorem/ZeroFreeRegion}{Proof of the zero-free region}}
\end{proof}

For a fixed modulus \(Q\), one has \(f_\chi(1)\ne0\) for every character.
By continuity the hypothesis near \(1\) holds for all sufficiently large
\(\mathcal L\). For the principal character \(f_\chi(1)\) is the residue at
the pole, and for a nonprincipal character the assertion is Dirichlet's
nonvanishing at \(1\). See also \cite[Chapter 11]{mv} for the fixed-modulus
zero-free region.
\leanlinks{\leanref{eventually_norm_logDeriv_entireL_le}}

\begin{lemma}\label{lem:fixed-modulus-pnt}
Fix an odd prime \(d\) and a positive integer \(q\) with \((d,q)=1\).
For every integer \(B\ge0\) and all sufficiently large \(y\), one has
\[
 |R_q(y)|\le(\log y)^{-B}.
\]
If \(z\ge y\ge1\), then \(|R_q(z)-R_q(y)|\le13(z-y)/y\).
\end{lemma}

\begin{proof}
Put \(\mathcal E_q(t)=\sum_{n\le t}\lambda_q(n)-t/\varphi(q)\).
The bounds \(\lambda_q\le2\Lambda\) and \(\sum_{n\le t}\Lambda(n)\le6t\)
give \(|\mathcal E_q(t)|\le13t\). Partial summation gives
\[
 R_q(y)=\int_1^y\mathcal E_q(t)\frac{dt}{t^2}-c_q.
\]
The bound for differences follows. For \(0\le u\le v\le u+1\), it gives
\(|R_q(e^v)-R_q(e^u)|\le13(e^{v-u}-1)\le26(v-u)\). Partitioning into intervals
of length at most \(1\) shows that \(U(v)=R_q(e^v)\) is Lipschitz on \(v\ge0\)
with constant \(26\).

Apply Lemma \ref{lem:uniform-log-derivative} simultaneously to the finite
families \(\psi,\psi\chi_d\) of fixed moduli \(q,qd\). Each difference in the
average in \eqref{eq:log-derivative} has modulus at most \(4\mathcal L^{25}\),
and \(1/|s|\le2\). Thus, for all sufficiently large \(\mathcal L\),
\(F(s)=G_q(s)/s\) is holomorphic on \(\RePart s\ge1-\mathcal L^{-26}\),
\(|s|\le e^\mathcal L\), with
\(|F(s)|\le2(4\mathcal L^{25}+1)\le\mathcal L^{26}\).
\leanlinks{\leanref{eventually_norm_logDerivMean_div_le}}

Take \(u\) large and put \(k=\lfloor u^{2B+4}\rfloor\), \(H=(k+1)/u\),
\(\mathcal L=(2B+5)\log u\), and \(\eta=\mathcal L^{-26}\).
Write
\[
 A=\frac{(-H)^k H}{k!}\mathcal T_q^{(k)}(1+H)
\]
for the normalized Taylor coefficient of
\(\mathcal T_q(s)=(F(s)-F(1))/(s-1)\) at \(1+H\).
Differentiation under the integral expresses \(A\) as the average of \(U(v)\)
against the density \(H^{k+1}v^k e^{-Hv}/k!\), which has mean \(u\) and variance
\(u^2/(k+1)\). For any \(\delta>0\),
\(|v-u|\le\delta+(v-u)^2/\delta\). Combining this inequality with the Lipschitz
bound and choosing \(\delta=u^{-B-1}\) gives \(|A-U(u)|\le52u^{-B-1}\).

The disk of center \(1+H\) and radius \(H+\eta\) lies in the holomorphy region
above. On its boundary \(|s-1|\ge\eta\), so
\(|\mathcal T_q(s)|\le 2\mathcal L^{26}/\eta\). Cauchy's estimate and
\(H/(H+\eta)\le e^{-\eta/(2H)}\) give
\[
 |A|\le4\mathcal L^{26}H\eta^{-1}
       \exp\!\left(-\frac{(k+1)\eta}{2H}\right)
      =4\mathcal L^{52}H e^{-u\eta/2}.
\]
For sufficiently large \(u\) we have \(\mathcal L^{52}\le u\),
\(4H\le u^{2B+5}\), and \(u^{3B+7}\le e^{\sqrt{u}/2}\), so
\(|A|\le u^{-B-1}\). Consequently \(|U(u)|\le53u^{-B-1}\le u^{-B}\).
The substitution \(u=\log y\) finishes the proof.
\leanlinks{\leanref{eventually_abs_le_inv_pow_of_transform}}
\end{proof}
